\documentclass[11pt]{article}
\usepackage[T1]{fontenc}
\usepackage[a4paper,margin=1in]{geometry}
\usepackage{amsmath,amssymb,amsthm}
\usepackage[hidelinks]{hyperref}
\usepackage{microtype}

\newtheorem{theorem}{Theorem}
\newtheorem{proposition}[theorem]{Proposition}
\newtheorem{lemma}[theorem]{Lemma}
\newtheorem*{sourcequestion}{Source question (Remark 5.9)}
\newtheorem*{answer}{Literature-implied answer}

\newcommand{\N}{\mathbb N}
\newcommand{\R}{\mathbb R}
\newcommand{\C}{\mathbb C}
\newcommand{\T}{\mathbb T}
\newcommand{\Lin}{\mathcal L}

\setlength{\parindent}{0pt}
\setlength{\parskip}{0.55em}

\title{A Literature-Implied Counterexample to\\
$L^\infty$-BIBO/Admissibility Equivalence}
\author{Run \texttt{fa\_banach\_001} --- lane 7}
\date{11 August 2026}

\begin{document}
\maketitle

\begin{center}
\fbox{\parbox{0.9\textwidth}{
\textbf{Status: full negative literature-implied answer.}
Kato's diagonal semigroup on $c_0$, recorded in Example~2.3 of
arXiv:2008.00459, makes the identity-output system $L^\infty$-BIBO stable
although its control operator is not $L^\infty$-admissible.  This answers
Remark~5.9 of arXiv:2408.06313 negatively.  The construction itself is
known literature; the relation to the later question is the identification
recorded here.
}}
\end{center}

\section{Question and answer}

Let $A$ generate a $C_0$-semigroup on a Banach space $X$, let
$B\in\Lin(U,X_{-1})$, and consider the identity-output system node
\[
 \Sigma_B=\Sigma\bigl(A,B,I,(\,\cdot\, I-A)^{-1}B\bigr).
\]
Schwenninger and Wierzba prove in Proposition~5.8 of \cite{source} that
$\Sigma_B$ is $C^\infty$-BIBO stable if and only if $B$ is infinite-time
$C$-control-admissible.  On arXiv PDF page~12 they ask:

\begin{sourcequestion}
Does the analogous equivalence hold for $L^\infty$ inputs?
\end{sourcequestion}

\begin{answer}
No.  There is an exponentially stable positive diagonal semigroup on $c_0$
and a positive control operator $B$ for which $\Sigma_B$ is
$L^\infty$-BIBO stable with constant one, while $B$ is not even finite-time
$L^\infty$-control-admissible.
\end{answer}

\section{The Kato system}

Let $X=U=c_0(\N)$ and define
\[
 (T(t)x)_n=e^{-nt}x_n,\qquad
 (Ax)_n=-n x_n,
\]
with
\[
 D(A)=\{x\in c_0:(n x_n)_{n\ge1}\in c_0\}.
\]
Then $(T(t))_{t\ge0}$ is a positive exponentially stable $C_0$-semigroup.
Using $0\in\rho(A)$ to norm the extrapolation space,
\[
 X_{-1}=\{z=(z_n):(z_n/n)\in c_0\},
 \qquad \|z\|_{-1}=\sup_n\frac{|z_n|}{n}.
\]
Set $B=-A_{-1}|_X$, so $(Bu)_n=n u_n$.  This is a bounded positive map
$c_0\to X_{-1}$.  Also
\[
 \bigl((\lambda I-A)^{-1}Bu\bigr)_n
 =\frac{n}{\lambda+n}u_n\in c_0,
\]
so $\Sigma_B$ is a legitimate system node with output space $c_0$.
Its combined observation/feedthrough operator is the identity: for every
admissible pair $(x,u)$,
\[
 (C\&D)\binom xu
 =x-(\lambda I-A)^{-1}Bu+(\lambda I-A)^{-1}Bu=x.
\]
Consequently its distributional output is the state itself.

For $u\in L^\infty_{\mathrm{loc}}([0,\infty);c_0)$ and zero initial state,
the extrapolation-space state has coordinates
\begin{equation}\label{eq:state}
 x_n(t)=n\int_0^t e^{-n(t-s)}u_n(s)\,ds.
\end{equation}

\section{$L^\infty$-BIBO stability}

The following elementary verification makes the system-node identification
with the maximal-regularity statement in \cite{support} explicit.

\begin{proposition}\label{prop:bibo}
For every $u\in L^\infty_{\mathrm{loc}}([0,\infty);c_0)$, the state in
\eqref{eq:state} has a $c_0$-valued strongly measurable representative and,
for every $\tau>0$,
\[
 \|x\|_{L^\infty([0,\tau];c_0)}
 \le \|u\|_{L^\infty([0,\tau];c_0)}.
\]
Thus $\Sigma_B$ is $L^\infty$-BIBO stable.
\end{proposition}

\begin{proof}
The coordinate estimate from \eqref{eq:state} gives, for almost every
$t\le\tau$,
\[
 |x_n(t)|\le (1-e^{-nt})
 \|u\|_{L^\infty([0,t];c_0)}
 \le \|u\|_{L^\infty([0,\tau];c_0)}.
\]
It remains to prove that $x_n(t)\to0$ for almost every $t$.  Extend $u$ by
zero to negative times and take a Lebesgue point $t$ of the Bochner function
$u$.  Writing $r=t-s$ gives
\[
 x_n(t)=(1-e^{-nt})u_n(t)
 +n\int_0^t e^{-nr}\bigl(u_n(t-r)-u_n(t)\bigr)\,dr.
\]
The first term tends to zero because $u(t)\in c_0$.  The absolute value of
the second is bounded by
\[
 n\int_0^t e^{-nr}\|u(t-r)-u(t)\|_{c_0}\,dr,
\]
which tends to zero because $(n e^{-nr})_{n\ge1}$ is a one-sided approximate
identity and $t$ is a Lebesgue point.  Hence $x(t)\in c_0$ for almost every
$t$.

Each coordinate is measurable.  On the full-measure set just obtained,
finite-coordinate truncations converge pointwise in $c_0$ norm to $x(t)$;
therefore $x$ has a strongly measurable $c_0$-valued representative.  The
coordinate estimate then proves the asserted essential-supremum bound.
\end{proof}

The supporting paper states the same regularity fact in its original
language: immediately after Example~2.3, it records that Kato's semigroup has
$L^\infty$-maximal regularity.  Since $B=-A_{-1}$, that assertion is exactly
the a.e. $c_0$-valued boundedness of the state above.

\section{Failure of endpoint admissibility}
\enlargethispage{7\baselineskip}

Fix $\tau>0$.  Example~2.3 of \cite{support}, on arXiv PDF pages~7--8,
uses the strongly measurable input $u\in L^\infty([0,\tau];c_0)$ given by
\[
 u_n(s)=
 \begin{cases}
 1,&s\in[\tau-1/n,\,\tau-1/(2n)]\text{ and }1/n<\tau,\\
 0,&\text{otherwise}.
 \end{cases}
\]
For each fixed $s$, only finitely many coordinates are nonzero, so
$u(s)\in c_0$, and $\|u\|_{L^\infty}=1$.  But at the endpoint,
\[
 x_n(\tau)
 =\int_{1/(2n)}^{1/n} n e^{-nr}\,dr
 =e^{-1/2}-e^{-1}>\frac15
\]
for every sufficiently large $n$.  Thus $x(\tau)\notin c_0$.  The endpoint
input map
\[
 \Phi_\tau u=\int_0^\tau T_{-1}(\tau-s)Bu(s)\,ds
\]
does not have range in $X$, so $B$ is not $L^\infty$-control-admissible.
Together with Proposition~\ref{prop:bibo}, this is the desired counterexample.

The mechanism is sharp: BIBO stability only sees the output as an almost-
everywhere equivalence class, whereas control admissibility requires the state
to lie in $X$ at the specified endpoint.  Kato's input creates a single
non-$c_0$ endpoint without violating the a.e. BIBO estimate.

\section{Reflexive-state boundary}

The endpoint defect cannot occur on reflexive state spaces.

\begin{proposition}\label{prop:reflexive}
If $X$ is reflexive, then for the identity-output system $\Sigma_B$,
$L^\infty$-BIBO stability is equivalent to infinite-time
$L^\infty$-control-admissibility.
\end{proposition}

\begin{proof}
Only the forward implication needs comment.  Fix $\tau>0$ and
$u\in L^\infty([0,\tau];U)$, extended by zero after $\tau$.  Let $x$ be the
zero-state trajectory.  BIBO stability gives $x(t)\in X$ for almost every
$t<\tau$ and $\|x(t)\|_X\le c\|u\|_\infty$.  Choose $t_k\uparrow\tau$ from
this full-measure set.  Reflexivity gives a subsequence converging weakly in
$X$ to some $z\in X$.  The canonical embedding $X\hookrightarrow X_{-1}$ is
continuous, while the generalized state is continuous in $X_{-1}$; hence the
same subsequence converges in $X_{-1}$ norm to $x(\tau)$.  Uniqueness of weak
limits in $X_{-1}$ gives $x(\tau)=z\in X$, with
$\|x(\tau)\|_X\le c\|u\|_\infty$.  The constant is independent of $\tau$,
which is infinite-time admissibility.  The converse follows directly from
the state formula and the identity output.
\end{proof}

This addendum is not needed for the negative answer, but it isolates the role
of the nonreflexive space $c_0$.

\section{Provenance and remaining scope}

The semigroup, the operator $B=-A_{-1}$, the explicit bad input, the
$L^\infty$-maximal-regularity property, and the failure of
$L^\infty$-admissibility are all contained in Example~2.3 and the following
paragraph of \cite{support}.  The present packet makes the direct translation
to Remark~5.9 of \cite{source}; no new counterexample is claimed.

This does not settle the broader Outlook question in \cite{source} asking
whether $C^\infty$-, $C$-, and $L^\infty$-BIBO stability are equivalent for
arbitrary system nodes.  In the Kato identity-output example all three BIBO
properties hold; what fails is endpoint $L^\infty$ control admissibility.

\small
\begin{thebibliography}{9}
\bibitem{source}
F. L. Schwenninger and A. A. Wierzba,
\emph{A dual notion to BIBO stability},
J. Math. Anal. Appl. \textbf{545} (2025), 129156;
arXiv:2408.06313.

\bibitem{support}
B. Jacob, F. L. Schwenninger, and J. Wintermayr,
\emph{A refinement of Baillon's theorem on maximal regularity},
Stud. Math. \textbf{263} (2022), 141--158;
arXiv:2008.00459.
\end{thebibliography}

\end{document}
